\documentclass[aspectratio=169,11pt]{beamer}
\usepackage{fontspec}
\setsansfont{DejaVu Sans}
\usepackage{polyglossia}
\setdefaultlanguage{russian}
\usepackage{booktabs}
\definecolor{ink}{HTML}{17324D}
\definecolor{accent}{HTML}{277D77}
\setbeamercolor{normal text}{fg=ink,bg=white}
\setbeamertemplate{navigation symbols}{}
\setbeamertemplate{footline}{}
\setbeamersize{text margin left=.55cm,text margin right=.55cm}
\begin{document}
\begin{frame}[t,plain]
{\fontsize{14}{17}\selectfont\bfseries Выбор роутера при дисбалансе обучения}\par
\vspace{.12cm}
{\fontsize{8}{10}\selectfont Anastasia Shlopak \hfill Руководитель: Archil Maysuradze, PhD}\par
\vspace{.15cm}
{\small Все 8 обучаемых методов · 21 пара · 7 составов · 20 повторений}\par
\vspace{.08cm}\centering\fontsize{8}{9}\selectfont
\renewcommand{\arraystretch}{1.05}
\begin{tabular}{lrrrrrr}
\toprule
 & \multicolumn{3}{c}{TF--IDF $\to$ BERT-tiny} & \multicolumn{3}{c}{TF--IDF $\to$ MiniLM} \\
Метод & Natural & Rescue8 & Harm8 & Natural & Rescue8 & Harm8 \\
\midrule
Weak-error & 87.87 & 87.89 & 87.67 & 91.26 & 91.30 & 91.25 \\
Direct & 87.61 & 88.22 & 86.38 & 91.44 & 91.57 & 86.56 \\
Marginal & 87.58 & 87.92 & 86.41 & 91.36 & 91.31 & 86.50 \\
Events & 87.50 & 87.93 & 86.40 & 91.45 & 91.35 & 86.50 \\
Ternary & 87.37 & 87.95 & 86.43 & 91.46 & 91.33 & 86.61 \\
Joint & 87.42 & 87.93 & 86.43 & 91.43 & 91.32 & 86.61 \\
Hurdle & 87.82 & 87.92 & 86.43 & 91.34 & 91.38 & 86.61 \\
Soft-gain & 87.59 & 88.04 & 86.38 & 91.37 & 91.33 & 86.52 \\
Confidence & 87.24 & 87.24 & 87.24 & 90.68 & 90.68 & 90.68 \\
\bottomrule
\end{tabular}
\par
\vspace{.08cm}
{\fontsize{8}{10}\selectfont Средняя accuracy, \%; fit 1\,000; квота 20\%. Два примера из полного контроля.}\par
\vspace{.15cm}
{\small\color{accent}\bfseries Harm-дисбаланс разделяет gain-цели и weak-error/confidence.}\par
\vspace{.12cm}
{\fontsize{8}{10}\selectfont 56 взаимодействий после общей поправки; между семью gain-методами — 0 из 882.\newline
Различие чувствительности не означает надёжную смену победителя. Перенос ещё не проверен.}
\end{frame}
\end{document}
